%% Correction audit de veracite - 2026-09-20 : citations de rapport de jury sans
%% support textuel dans le corpus (rapports 2018-2025) remplacees par des verbatims
%% verifies ou supprimees. Detail ci-dessous.
%% l.319 (ex) : citation "jury 2025" reformulee (lois basiques / fonctions de transfert /
%%   redaction des demonstrations) INTROUVABLE telle quelle -> remplacee par le verbatim
%%   exact de RJ2025 l.484. l.308 et l.320 sont conformes : non touchees.
\expandafter\ifx\csname CAPETINCLUDEMODE\endcsname\relax
\documentclass[11pt,a4paper]{article}
% preamble.tex — packages et commandes communes a toutes les fiches CAPET SII IE
\input{glyphtounicode}
\pdfgentounicode=1
\usepackage{cmap}
\usepackage[T1]{fontenc}
\usepackage[utf8]{inputenc}
\usepackage[french]{babel}
% babel[french] active a \begin{document} des raccourcis de ponctuation (:;!?)
% qui cassent la syntaxe to[...] de circuitikz (erreur "+ or - expected") ;
% on les desactive, ce qui ne change que l'espacement fin avant ces signes.
\AtBeginDocument{\shorthandoff{:;!?}}
\usepackage[a4paper, top=1.6cm, bottom=1.6cm, left=1.4cm, right=1.4cm]{geometry}
\setlength{\headheight}{14.5pt}
\usepackage{amsmath, amssymb, amsfonts}
\usepackage{siunitx}
\sisetup{locale=FR, output-decimal-marker={,}, inter-unit-product=\ensuremath{{}\cdot{}},
         group-minimum-digits=4}
% Unites non fournies par siunitx, declarees ici une fois pour toutes (fiches + exercices).
% Les gardes \ifdefined evitent tout conflit avec une redeclaration locale restee dans un
% fichier : ne pas retirer les gardes, ne pas redeclarer ces unites dans les fichiers.
\DeclareSIUnit{\jour}{j}
\ifdefined\var\else\DeclareSIUnit{\var}{var}\fi   % puissance reactive
\ifdefined\vad\else\DeclareSIUnit{\vad}{VAD}\fi   % puissance deformante
\ifdefined\tr\else\DeclareSIUnit{\tr}{tr}\fi      % tour (tr/min)
\usepackage{array, booktabs, multirow, tabularx}
\usepackage{xcolor}
\usepackage[most, breakable]{tcolorbox}
\usepackage{multicol}
\usepackage{enumitem}
\usepackage{mathtools}
\usepackage{lmodern}
% Refonte visuelle 2026 (direction "moderne coloré") : Lato pour le texte,
% maths sans-serif assorties via newtxsf (charge apres amsmath/amssymb).
\usepackage[default]{lato}
\usepackage{newtxsf}
\usepackage{microtype}
\usepackage{fancyhdr}
\usepackage{tikz}
\usetikzlibrary{positioning,arrows.meta}
\usepackage{circuitikz}
\usepackage{lastpage}

% ── Couleurs par domaine (palette modernisée 2026) ────────────────────
\definecolor{colElec}{RGB}{29, 78, 216}        % Électrotechnique — bleu vif
\definecolor{colElecLight}{RGB}{219, 234, 254} % fond bleu clair
\definecolor{colAuto}{RGB}{21, 128, 61}        % Automatique — vert
\definecolor{colMeca}{RGB}{194, 65, 12}        % Mécanique — orange
\definecolor{colTherm}{RGB}{185, 28, 28}       % Thermique — rouge
\definecolor{colQT}{RGB}{232, 244, 233}        % (hérité, plus utilisé par questiontype)
\definecolor{colJury}{RGB}{255, 248, 225}      % fond jury (crème)
\definecolor{colJuryFrame}{RGB}{202, 138, 4}   % bordure/bandeau jury (ambre)
\definecolor{colGrayLight}{RGB}{245, 245, 247} % fond gris neutre
\definecolor{colSeparator}{RGB}{180, 180, 180} % gris séparateur
\definecolor{colExo}{RGB}{8, 145, 178}          % Banque d'exercices — cyan (distinct des fiches)
% ── Encre et gris de la refonte ───────────────────────────────────────
\definecolor{colInk}{RGB}{23, 32, 42}          % texte courant (quasi-noir)
\definecolor{colSlate}{RGB}{83, 98, 118}       % gris ardoise (méta, en-têtes)
\definecolor{colRule}{RGB}{203, 213, 225}      % filets clairs
\definecolor{colJuryBg}{RGB}{255, 248, 225}    % fond crème jury
\definecolor{colPresence}{RGB}{83, 98, 118}    % bloc 5 (présence dans les sujets)
\AtBeginDocument{\color{colInk}}

% ── Style encadré exercice + corrigé (banque d'exercices) ─────────────
\tcbset{
  exoBox/.style={
    breakable,
    enhanced,
    colback=colExo!8!white,
    colframe=colExo,
    colbacktitle=colExo,
    coltitle=white,
    fonttitle=\bfseries\small,
    left=4pt, right=4pt, top=3pt, bottom=3pt,
    before skip=5pt, after skip=5pt,
  }
}

% ── Bascule enonce / corrige (banque d'exercices) ─────────────────────
% \CAPETCORRIGE est defini en ligne de commande (pdflatex "\def\CAPETCORRIGE{1}\input{...}")
% pour compiler la variante corriges d'un meme fichier source ; par defaut (non defini),
% la compilation produit la variante enonces. \CAPETSHOW{enonce}{corrige} choisit la branche.
%
% TROIS MODES (refonte 2026-09-20, manuel de revision) :
%   - aucun drapeau   -> imprime l'enonce            (variante enonces, INCHANGE)
%   - \CAPETCORRIGE   -> imprime le corrige          (variante corriges, INCHANGE)
%   - \CAPETMANUEL    -> imprime l'enonce ET met le corrige de cote pour l'annexe
%
% Le mode manuel stocke chaque corrige dans une macro (\csgdef), sans ecriture
% fichier ni re-tokenisation : le contenu (circuitikz, accents, %) traverse intact.
% L'annexe les rejoue ensuite dans l'ordre via \CAPETANNEXE.
\usepackage{etoolbox}

\newcounter{CAPETexo}          % numero continu d'exercice dans le manuel
\setcounter{CAPETexo}{0}
\newcounter{CAPETannexe}       % compteur de relecture en annexe

% Titre courant de la notion en cours (pose par \CAPETNOTION), sert d'intertitre annexe
\newcommand{\CAPETnotioncode}{}
\newcommand{\CAPETnotiontitre}{}

% \CAPETNOTION{code}{titre} : ouvre une notion dans le manuel (memorise le libelle,
% et marque le rang du premier exercice de la notion pour l'intertitre d'annexe).
\newcommand{\CAPETNOTION}[2]{%
  \renewcommand{\CAPETnotioncode}{#1}%
  \renewcommand{\CAPETnotiontitre}{#2}%
  \csgdef{CAPETnotionstart\the\numexpr\value{CAPETexo}+1\relax}{#1 -- #2}%
}

% ── \degree robuste ───────────────────────────────────────────────────
% Certains fichiers d'exercices utilisent \degree. Ni le preambule ni siunitx
% ne le definissent, et \providecommand local ne suffit pas : sous hyperref,
% \degree existe deja comme commande factice (\ERROR), donc \providecommand
% s'abstient et l'appel echoue a la compilation du manuel (erreur invisible
% en compilation isolee). On le (re)definit donc inconditionnellement ici,
% apres le chargement des packages, pour tout le corpus.
\AtBeginDocument{\providecommand{\degree}{}\renewcommand{\degree}{\ensuremath{^\circ}}}

% Libelle decoratif : les fichiers d'exercices utilisent aussi \CAPETSHOW dans
% leur \fancyhead et leur bloc de titre, pour afficher "(enonces)" ou "(corriges)".
% Ces appels-la ne sont PAS des exercices : en mode manuel ils ne doivent ni
% consommer un numero, ni alimenter l'annexe. On les reconnait a leur premier
% argument, strictement egal au libelle ci-dessous.
\newcommand{\CAPETLIBELLEENONCES}{(\'enonc\'es)}

\ifdefined\CAPETMANUEL
  % --- Mode manuel : enonce en place, corrige differe vers l'annexe ---
  \newcommand{\CAPETSHOW}[2]{%
    \def\CAPETargun{#1}%
    \ifx\CAPETargun\CAPETLIBELLEENONCES
      #1%            appel decoratif : rendu tel quel, hors numerotation
    \else
      \CAPETSHOWEXO{#1}{#2}%
    \fi
  }
  \newcommand{\CAPETSHOWEXO}[2]{%
    \stepcounter{CAPETexo}%
    \csgdef{CAPETcorrige\theCAPETexo}{#2}%
    \edef\CAPETtmpnum{\theCAPETexo}%
    #1%
    \par\vspace{2pt}%
    {\footnotesize\itshape\color{colExo}%
     \hyperref[capet:cor:\CAPETtmpnum]{Corrige de l'exercice \CAPETtmpnum\ $\rightarrow$ p.~\pageref{capet:cor:\CAPETtmpnum}}%
     \label{capet:eno:\CAPETtmpnum}}%
  }
  % \CAPETANNEXE : rejoue tous les corriges captures, dans l'ordre, avec intertitres
  \newcommand{\CAPETANNEXE}{%
    \setcounter{CAPETannexe}{0}%
    \loop\ifnum\value{CAPETannexe}<\value{CAPETexo}%
      \stepcounter{CAPETannexe}%
      \edef\CAPETcurnum{\theCAPETannexe}%
      \ifcsname CAPETnotionstart\CAPETcurnum\endcsname
        \par\medskip
        {\normalsize\bfseries\color{colExo}\csuse{CAPETnotionstart\CAPETcurnum}}%
        \par\nobreak\smallskip
      \fi
      \begin{tcolorbox}[exoBox, breakable,
        title={Corrige de l'exercice \CAPETcurnum\
               \hfill\normalfont\footnotesize
               \hyperref[capet:eno:\CAPETcurnum]{enonce p.~\pageref{capet:eno:\CAPETcurnum}}}]%
      \label{capet:cor:\CAPETcurnum}%
      \csuse{CAPETcorrige\CAPETcurnum}%
      \end{tcolorbox}%
    \repeat
  }
\else
  % --- Modes historiques, comportement strictement inchange ---
  \newcommand{\CAPETSHOW}[2]{\ifdefined\CAPETCORRIGE #2\else #1\fi}
  \newcommand{\CAPETANNEXE}{}
\fi

% ── Difficulté 1 à 3 (étoiles pleines, amssymb) ───────────────────────
\newcommand{\difficulte}[1]{%
  \ifcase#1\relax
  \or $\bigstar$\;niveau 1/3%
  \or $\bigstar\bigstar$\;niveau 2/3%
  \or $\bigstar\bigstar\bigstar$\;niveau 3/3%
  \fi
}

% ── Style encadré extrait de document technique (banque d'exercices) ──
\tcbset{
  docTech/.style={
    breakable,
    enhanced,
    colback=colGrayLight,
    colframe=colSeparator,
    fonttitle=\bfseries\small\itshape,
    coltitle=black,
    colbacktitle=colGrayLight,
    left=4pt, right=4pt, top=3pt, bottom=3pt,
    before skip=4pt, after skip=4pt,
  }
}

% ── Style tcolorbox de section principale ────────────────────────────
% Refonte 2026 : bandeau titre coloré arrondi + fond très pâle, sans cadre.
\tcbset{
  ficheSection/.style={
    breakable, enhanced jigsaw,
    frame hidden, boxrule=0pt, arc=1.8mm,
    colback=#1!4!white,
    colbacktitle=#1,
    coltitle=white,
    fonttitle=\sffamily\bfseries,
    left=7pt, right=7pt, top=5pt, bottom=5pt,
    toptitle=3pt, bottomtitle=3pt,
    before skip=10pt, after skip=8pt,
  }
}

% ── Style encadré question type ──────────────────────────────────────
% Carte blanche, bandeau sombre, liseré très clair.
\tcbset{
  questiontype/.style={
    breakable, enhanced jigsaw,
    colback=white,
    frame hidden, boxrule=0pt, arc=1.8mm,
    colbacktitle=colElec!78!black,
    coltitle=white,
    fonttitle=\sffamily\bfseries\small,
    left=7pt, right=7pt, top=4pt, bottom=4pt,
    toptitle=2.5pt, bottomtitle=2.5pt,
    before skip=8pt, after skip=8pt,
    borderline={0.6pt}{0pt}{colElec!12!white},
  }
}

% ── Style encadré jury ───────────────────────────────────────────────
% Bandeau ambre, fond crème, sans cadre.
\tcbset{
  jury/.style={
    breakable, enhanced jigsaw,
    colback=colJuryBg,
    frame hidden, boxrule=0pt, arc=1.8mm,
    colbacktitle=colJuryFrame,
    coltitle=white,
    fonttitle=\sffamily\bfseries,
    left=7pt, right=7pt, top=4pt, bottom=4pt,
    toptitle=3pt, bottomtitle=3pt,
    before skip=10pt, after skip=8pt,
  }
}

% ── En-tête / pied de page communs (refonte 2026) ─────────────────────
% #1 code fiche, #2 titre court, #3 couleur domaine
\newcommand{\ficheHeaders}[3]{%
  \fancyhead[L]{\small\sffamily{\bfseries\textcolor{#3}{#1}}\hspace{0.7em}\textcolor{colSlate}{#2}}%
  \fancyhead[R]{\small\sffamily\textcolor{colSlate}{CAPET SII IE\, --- \,Session 2026}}%
  \fancyfoot[C]{\small\sffamily\textcolor{colSlate}{\thepage\,/\,\pageref{LastPage}}}%
  \fancyhead[C]{}\fancyfoot[L]{}\fancyfoot[R]{}%
}

% ── Bloc en-tête de fiche (BLOC 0, refonte 2026) ──────────────────────
% #1 code  #2 titre  #3 couleur domaine  #4 domaine  #5 ligne présence  #6 sous-thèmes
\newcommand{\ficheEntete}[6]{%
  \begin{tcolorbox}[enhanced, breakable,
    frame hidden, boxrule=0pt, arc=2.5mm,
    interior style={left color=#3, right color=#3!75!black},
    left=9pt, right=9pt, top=7pt, bottom=7pt,
    before skip=0pt, after skip=6pt]
    {\color{white}\sffamily\bfseries\LARGE #1\quad #2}\\[4pt]
    {\color{#3!22!white}\sffamily\small \MakeUppercase{#4}\quad\textbar\quad #5}\\[1pt]
    {\color{#3!35!white}\sffamily\footnotesize Sous-thèmes : #6}
  \end{tcolorbox}%
}

% ── Mise en page ─────────────────────────────────────────────────────
\pagestyle{fancy}
\fancyhf{}
\renewcommand{\headrulewidth}{0.5pt}
\renewcommand{\footrulewidth}{0pt}
\AtBeginDocument{%
  \renewcommand{\headrule}{{\color{colRule}\hrule width\headwidth height 0.5pt}}%
}

\setlist[itemize]{leftmargin=1.2em, itemsep=1pt, topsep=2pt, parsep=0pt}
\setlist[enumerate]{leftmargin=1.2em, itemsep=3pt, topsep=2pt, parsep=0pt}

\setlength{\parskip}{2pt}
\setlength{\columnsep}{10pt}

\def\CAPETBEGINFILE{\begin{document}}
\def\CAPETENDFILE{\end{document}}
\else
\def\CAPETBEGINFILE{}
\def\CAPETENDFILE{}
\fi

\ficheHeaders{T21}{Fonctions de transfert}{colElec}

\CAPETBEGINFILE

% -----------------------------------------------------------------------
%  BLOC EN-TETE
% -----------------------------------------------------------------------

\ficheEntete{T21}{Fonctions de transfert - Diagramme de Bode}{colElec}%
  {électronique / Automatique - Analyse fréquentielle}%
  {\textbf{Présence :} \textbf{8/13} sujets (2013-2026) - thème \emph{fondamental}}%
  {transformée de Laplace - FdT 1\textsuperscript{er}/2\textsuperscript{e} ordre - identification graphique Bode - marges de phase et de gain - correcteur proportionnel - erreur statique}

% -----------------------------------------------------------------------
%  1. RAPPELS DE COURS
% -----------------------------------------------------------------------

\begin{tcolorbox}[
  ficheSection=colElec,
  title={1. Rappels de cours},
]

% -- 1.1 Transformee de Laplace ----------------------------------------
\textbf{1.1 Transformée de Laplace - rappels essentiels}

\smallskip
Conditions de Heaviside (conditions initiales nulles) :
\[
  \mathcal{L}\!\left\{\frac{dx}{dt}\right\} = p\,X(p)
  \qquad
  \mathcal{L}\!\left\{\int x\,dt\right\} = \frac{X(p)}{p}
  \qquad
  H(p) = \frac{S(p)}{E(p)}
\]

\medskip\hrule\vspace{4pt}

% -- 1.2 FdT 1er ordre ------------------------------------------------
\textbf{1.2 Système du 1\textsuperscript{er} ordre - forme canonique}

\[
  H(p) = \frac{K}{1 + \tau\,p}
  \qquad K : \text{gain statique} \qquad \tau : \text{constante de temps}
\]

Réponse indicielle à l'échelon $E_0$ : $s(t) = K\,E_0\,(1 - e^{-t/\tau})$ \quad;\quad $t_{r5\%} \approx 3\tau$

\textbf{Identification sur le Bode :}
\begin{itemize}
  \item Plateau basse fréquence $= 20\log K$ dB
  \item Pulsation de coupure = intersection des asymptotes 0\,dB/dec et $-20$\,dB/dec
\end{itemize}

\medskip\hrule\vspace{4pt}

% -- 1.3 FdT 2e ordre -------------------------------------------------
\textbf{1.3 Système du 2\textsuperscript{e} ordre - forme canonique}

\[
  H(p) = \frac{K}{1 + \dfrac{2m}{\omega_0}\,p + \dfrac{p^2}{\omega_0^2}}
\]
\[
  D\% = 100\,e^{-m\pi/\sqrt{1-m^2}} \quad (m<1) \qquad;\qquad t_{r5\%}\cdot\omega_0 = f(m) \text{ (abaque)}
\]

\medskip\hrule\vspace{4pt}

% -- 1.4 Diagramme de Bode asymptotique --------------------------------
\textbf{1.4 Bode asymptotique - règles de tracé}

\smallskip
\begin{center}
\renewcommand{\arraystretch}{1.3}
\begin{tabular}{llll}
\toprule
\textbf{Facteur} & \textbf{Gain asymptotique} & \textbf{Phase} & \textbf{Pente} \\
\midrule
Gain pur $K$ & $20\log K$ (constant) & $0^\circ$ & 0 \\
Intégrateur $1/p$ & $0\,\mathrm{dB}$ à $\omega=1$ & $-90^\circ$ & $-20\,\mathrm{dB/dec}$ \\
1\textsuperscript{er} ordre $(1+\tau p)^{-1}$ & $0$ puis $-20\,\mathrm{dB/dec}$ à $1/\tau$ & $0^\circ$ puis $-90^\circ$ & \\
2\textsuperscript{e} ordre & $0$ puis $-40\,\mathrm{dB/dec}$ à $\omega_0$ & $0^\circ$ puis $-180^\circ$ & \\
\bottomrule
\end{tabular}
\end{center}
Pour $m < 1/\sqrt{2} \approx 0{,}707$ : pic de résonance $G_{pic} = 20\log\dfrac{1}{2m\sqrt{1-m^2}}$.

\smallskip
\begin{center}
\begin{tikzpicture}[scale=1, font=\footnotesize]
  % --- Diagramme de gain (1er ordre) ---
  \node[below] at (2.6,-2.15) {\bfseries Gain};
  \draw[->] (0,0) -- (5.4,0) node[right] {$\log\omega$};
  \draw[->] (0,-1.8) -- (0,1.3) node[above] {$G_{dB}$};
  \draw[thick, colElec] (0,0.85) -- (2,0.85);
  \draw[thick, colElec] (2,0.85) -- (5,-1.4);
  \draw[dashed] (2,0.85) -- (2,-1.8);
  \node[below] at (2,-1.8) {$\omega_c$};
  \node[above] at (0.9,0.85) {$20\log K$};
  \node[colElec] at (4.1,-0.55) {\rotatebox{-27}{$-20\,\mathrm{dB/dec}$}};

  % --- Diagramme de phase (1er ordre) ---
  \begin{scope}[xshift=7.4cm]
    \node[below] at (2.5,-2.15) {\bfseries Phase};
    \draw[->] (0,0.3) -- (5.4,0.3) node[right] {$\log\omega$};
    \draw[->] (0,-1.8) -- (0,0.6) node[above] {$\varphi$};
    \draw[thick, colElec] (0,0.3) .. controls (1.6,0.3) and (2.4,-1.35) .. (4,-1.35);
    \draw[thick, colElec] (4,-1.35) -- (5,-1.35);
    \draw[dashed] (2,0.3) -- (2,-1.8);
    \node[below] at (2,-1.8) {$\omega_c$};
    \node[left] at (-0.1,-0.68) {$-45^\circ$};
    \node[left] at (-0.1,-1.35) {$-90^\circ$};
  \end{scope}
\end{tikzpicture}
\end{center}

\medskip\hrule\vspace{4pt}

% -- 1.5 Marges de stabilite -------------------------------------------
\textbf{1.5 Marges de phase et de gain (FTBO)}

\[
  M_\varphi = 180^\circ + \arg\!\bigl(FTBO(j\omega_u)\bigr) \quad \text{avec } |FTBO(j\omega_u)| = 1 \quad (0\,\mathrm{dB})
\]
\[
  M_G = -20\log|FTBO(j\omega_\pi)| \quad \text{avec } \arg(FTBO(j\omega_\pi)) = -180^\circ
\]

Critères habituels : $M_\varphi \geq 45^\circ$ (2023, 2026) ; $M_\varphi \geq 90^\circ$ (2024, 2025) ; $M_G \geq 10\,\mathrm{dB}$ (2026).

\textbf{Procédure graphique :}
\begin{enumerate}[leftmargin=1.5em]
  \item Lire $\omega_u$ sur la courbe de gain (croisement 0\,dB).
  \item Lire la phase à $\omega_u$ : $M_\varphi = 180^\circ + \varphi(\omega_u)$.
  \item Lire $\omega_\pi$ sur la courbe de phase (croisement $-180^\circ$).
  \item Lire le gain à $\omega_\pi$ : $M_G = -G_{dB}(\omega_\pi)$.
\end{enumerate}

\medskip\hrule\vspace{4pt}

% -- 1.6 Correcteur proportionnel --------------------------------------
\textbf{1.6 Correcteur proportionnel $C(p) = K_C$}

Ajouter $K_C$ décale le gain de $+20\log K_C$ dB \textbf{sans modifier la phase}.\\
Détermination de $K_C$ pour $M_\varphi$ imposée : lire $G_0$ dB à $\omega_u$ souhaité, puis $K_C = 10^{-G_0/20}$.

\end{tcolorbox}

% -----------------------------------------------------------------------
%  2. FORMULAIRE SYNTHETIQUE
% -----------------------------------------------------------------------

\begin{tcolorbox}[
  ficheSection=colElec,
  title={2. Formulaire synthétique},
]
\small
\begin{multicols}{2}
\renewcommand{\arraystretch}{1.45}
\begin{tabular}{@{}p{4.8cm}l@{}}
\toprule
\textbf{Formule} & \textbf{Usage} \\
\midrule
\multicolumn{2}{@{}l}{\footnotesize\textit{1\textsuperscript{er} ordre}} \\
$H(p) = K/(1+\tau p)$ & Forme canonique \\
$t_{r5\%} \approx 3\tau$ & Rapidité \\
$K = s(\infty)/E_0$ & Identification \\
$\tau$ : lu à $63{,}2\,\%$ de $K\,E_0$ & Identification \\
\midrule
\multicolumn{2}{@{}l}{\footnotesize\textit{2\textsuperscript{e} ordre}} \\
$H(p) = K/(1+2m\,p/\omega_0 + p^2/\omega_0^2)$ & Forme canonique \\
$D\% = 100\,e^{-m\pi/\sqrt{1-m^2}}$ & Dépassement \\
\bottomrule
\end{tabular}

\columnbreak

\renewcommand{\arraystretch}{1.45}
\begin{tabular}{@{}p{4.8cm}l@{}}
\toprule
\textbf{Formule} & \textbf{Usage} \\
\midrule
\multicolumn{2}{@{}l}{\footnotesize\textit{Marges de stabilité (FTBO)}} \\
$M_\varphi = 180^\circ + \varphi(\omega_u)$ & Marge de phase \\
$M_G = -G_{dB}(\omega_\pi)$ & Marge de gain \\
$M_\varphi \geq 45^\circ$ ; $M_G \geq 10\,\mathrm{dB}$ & CdC typique \\
\midrule
\multicolumn{2}{@{}l}{\footnotesize\textit{Correcteur et erreur}} \\
$K_C = 10^{-G_0/20}$ & Réglage $K_C$ \\
$\varepsilon_s = E/(1+K_{BO})$ & Erreur statique \\
$\varepsilon_s = 0$ si intégrateur BO & Classe 1 \\
\bottomrule
\end{tabular}
\end{multicols}
\normalsize
\end{tcolorbox}

% -----------------------------------------------------------------------
%  3. QUESTIONS TYPES
% -----------------------------------------------------------------------

\begin{tcolorbox}[
  ficheSection=colElec,
  title={3. Questions types - extraites des sujets réels},
]

% -- QT1 -------------------------------------------------------------

\begin{tcolorbox}[questiontype,
  title={QT1 - Identification graphique de $H_M(p)$ sur le Bode \hfill\normalfont\footnotesize E2 disc. 2023 (D4)}
]
\textbf{Énoncé.}
Le diagramme de Bode de $H_M(p)$ montre un plateau BF à $+26{,}4\,\mathrm{dB}$ et deux pulsations de coupure $\omega_a$ et $\omega_b$. Déterminer $K_M$, $\omega_a$, $\omega_b$ par lecture graphique. Proposer une expression approchée si $\tau_a \gg \tau_b$.

\tcbline

\textbf{Corrigé.}

\textbf{Étape 1 — Gain plateau BF.} $20\log K_M = 26{,}4\,\mathrm{dB} \Rightarrow K_M = 10^{26{,}4/20} \approx \mathbf{20{,}9}$.

\textbf{Étape 2 — Pulsations de coupure.} $\omega_a$ et $\omega_b$ se lisent aux intersections des asymptotes (0 dB/dec avec $-20$ dB/dec, puis $-20$ avec $-40$ dB/dec).

\textbf{Étape 3 — Modèle réduit.} Si $\tau_a \gg \tau_b$, le pôle en $\omega_a$ est dominant : $H_M(p) \approx K_M / (1 + \tau_a p)$.

\textbf{Conclusion.} Le modèle réduit du premier ordre reproduit fidèlement la réponse temporelle tant que la question ne porte pas sur le comportement haute fréquence : il satisfait l'exigence de simplification demandée par l'énoncé.
\end{tcolorbox}

\vspace{4pt}

% -- QT2 -------------------------------------------------------------

\begin{tcolorbox}[questiontype,
  title={QT2 - Marge de phase et réglage $K_C$ \hfill\normalfont\footnotesize E2 disc. 2023 (D7-D8)}
]
\textbf{Énoncé.}
FTBO tracée pour $K_C = 1$. La phase vaut $-135^\circ$ à $\omega_u$ (croisement 0\,dB). Calculer la marge de phase. Déterminer $K_C$ pour obtenir $M_\varphi = 45^\circ$.

\tcbline

\textbf{Corrigé.}

\textbf{Étape 1 — Marge de phase courante.} $M_\varphi = 180^\circ + (-135^\circ) = \mathbf{45^\circ}$ : la condition est déjà satisfaite pour $K_C = 1$.

\textbf{Étape 2 — Réglage si le gabarit change} \emph{[corrigé reconstitué : valeur $G_0$ illustrative, non issue du sujet]}\textbf{.} Si le gain lu à la pulsation souhaitée $\omega_u^*$ vaut $G_0 = +12\,\si{\decibel}$ :
\[
  K_C = 10^{-12/20} = 10^{-0{,}6} \approx \mathbf{0{,}25}
\]
Le décalage du gain de $20\log K_C = -12\,\si{\decibel}$ ramène le croisement $0\,\si{\decibel}$ à $\omega_u^*$ avec $M_\varphi = 45^\circ$.

\textbf{Conclusion.} Avec $K_C \approx 0{,}25$, la marge de phase obtenue respecte le critère de stabilité imposé ($M_\varphi \geq 45^\circ$) : le réglage est validé.
\end{tcolorbox}

\vspace{4pt}

% -- QT3 -------------------------------------------------------------

\begin{tcolorbox}[questiontype,
  title={QT3 - FTBF du 2\textsuperscript{e} ordre - dépassement nul \hfill\normalfont\footnotesize E1 2026 (C7-C8)}
]
\textbf{Énoncé.}
Montrer que la FTBF d'une boucle d'asservissement se met sous la forme du 2\textsuperscript{e} ordre. Exprimer $\omega_0$ et $m$ en fonction du gain $K_C$. Déduire la valeur de $K_C$ assurant $D\% = 0$.

\tcbline

\textbf{Corrigé.}

\textbf{Étape 1 — Identification canonique.} On identifie la FTBF sous forme canonique du second ordre : $\omega_0$ et $m$ s'expriment en fonction de $K_C$ selon le système étudié ($m$ croissant ou décroissant avec $K_C$ suivant les cas).

\textbf{Étape 2 — Condition de dépassement nul.} $D\% = 0 \Leftrightarrow m \geq 1$ (régime apériodique ou critique).

\textbf{Étape 3 — Application numérique.} Pour $m = 1$, on obtient $K_C = K_{C,lim}$, lu graphiquement sur l'abaque ou calculé analytiquement.

\textbf{Conclusion.} Pour $K_C = K_{C,lim}$, le système fonctionne à la limite de l'amortissement critique : le dépassement est annulé sans dégrader excessivement la rapidité, ce qui satisfait le cahier des charges — à condition de vérifier aussi que $M_G$ reste correct.
\end{tcolorbox}

\end{tcolorbox}

% -----------------------------------------------------------------------
%  4. ATTENTES DU JURY
% -----------------------------------------------------------------------

\begin{tcolorbox}[
  jury,
  title={4. Attentes du jury - Erreurs récurrentes et points de vigilance},
]

\begin{enumerate}[label=\textbf{\arabic*.}, leftmargin=*, itemsep=5pt]

  \item \textbf{Mettre la FdT sous forme canonique.} Le jury sanctionne l'absence de mise sous forme avec $K$, $\tau$ (ou $m$, $\omega_0$) explicitement identifiés.

  \item \textbf{Lire les marges graphiquement.} Ne pas chercher à les calculer analytiquement ; le document réponse fournit le diagramme de Bode : s'en servir.

  \item \textbf{Distinguer FTBO et FTBF.} Les marges se calculent sur la \textbf{boucle ouverte} ; la stabilité s'apprécie sur la \textbf{boucle fermée}.

  \item \textbf{Signe de la marge de phase.} $M_\varphi = 180^\circ + \varphi(\omega_u)$. Si $\varphi = -150^\circ$ alors $M_\varphi = 30^\circ$ (positif : stable).

  \item \textbf{Identification depuis l'essai indiciel.} Lire $\tau$ à $63{,}2\,\%$ (1\textsuperscript{er} ordre). Pour le 2\textsuperscript{e} ordre : lire le dépassement et la pseudo-période pour obtenir $m$ et $\omega_0$.

  \item \textbf{Thème le plus discriminant.} Jury 2024 : traité par 79\,\% des candidats, correctement par seulement 31\,\% — prévoir un entraînement intensif sur la lecture de Bode.

\end{enumerate}

\medskip
\begin{tcolorbox}[
  colback=colJuryFrame!12!colJury,
  colframe=colJuryFrame,
  left=4pt, right=4pt, top=2pt, bottom=2pt,
  fontupper=\small\itshape
]
Rapport jury 2025 : «~De nombreux candidats ne maitrisent pas les lois basiques de l'électricité (la loi des nœuds, la loi des mailles, les calculs de puissance ou encore la loi d'Ohm).~»\\
Rapport jury 2023 : «~L'automatique est un champ disciplinaire important des SII auquel les candidats ne peuvent se soustraire. Le jury invite à revoir les concepts de base.~»
\end{tcolorbox}

\end{tcolorbox}

% -----------------------------------------------------------------------
%  PIED - PRESENCE DANS LES SUJETS
% -----------------------------------------------------------------------

\begin{tcolorbox}[
  ficheSection=colPresence,
  title={Présence dans les sujets analysés (2013-2026)},
]
\begin{center}
\footnotesize\renewcommand{\arraystretch}{1.15}
\begin{tabular}{lccccccccccccc}
\toprule
\textbf{Année} & 13 & 14 & 15 & 16 & 17 & 18 & 19 & 20 & 21 & 22 & 23 & 24 & 25 \\
\midrule
E1 (commune)   & - & - & - & $\checkmark$ & - & - & $\checkmark$ & - & $\checkmark$ & - & $\checkmark$ & - & $\checkmark$ \\
E2 discipl. IE & - & $\checkmark$ & - & - & - & - & - & $\checkmark$ & - & - & $\checkmark$ & $\checkmark$ & $\checkmark$ \\
E2 appl. comm. & - & - & - & - & - & - & - & - & - & - & - & $\checkmark$ & - \\
\bottomrule
\end{tabular}
\end{center}
\smallskip
\footnotesize\textbf{Sous-thèmes fréquents :}
identification graphique FdT (Bode + essai indiciel) - mise sous forme canonique 1\textsuperscript{er}/2\textsuperscript{e} ordre - marges de phase et de gain - correcteur proportionnel $K_C$ - erreur statique (classe 0 vs 1)

\medskip
\textbf{Ordres de grandeur (repères rapides).}
\begin{itemize}[leftmargin=1.3em]
  \item Réseau BT FR : tension phase-neutre typique \SI{230}{\volt}, tension phase-phase typique \SI{400}{\volt}, fréquence \SI{50}{\hertz}.
  \item Courants d'étude fréquents en sujet : de quelques \si{\ampere} à quelques dizaines de \si{\ampere}.
  \item En conversion d'énergie, un rendement attendu reste généralement entre 0{,}8 et 0{,}98 selon la technologie.
\end{itemize}

\medskip
\textbf{Synthèse express (à refaire sans support).}
\begin{enumerate}[leftmargin=1.4em]
  \item Poser les hypothèses du modèle et les conventions de signe.
  \item Écrire 2 à 4 relations clés avec unités.
  \item Réaliser un calcul numérique complet et vérifier l'ordre de grandeur.
  \item Justifier le résultat par un argument technique de niveau CAPET.
\end{enumerate}

\end{tcolorbox}

\CAPETENDFILE
